\startslide
\commonframe
\mbox{}

\begin{tikzpicture}[remember picture,overlay]
  \slidetitle{Цель и задачи}

  \node[
    anchor=north west,
    fill=softbg,
    rounded corners=10pt,
    text width=5.9cm,
    inner xsep=0.35cm,
    inner ysep=0.28cm,
    align=left
  ] at ([xshift=1.00cm,yshift=-1.95cm]current page.north west) {%
    {\fontsize{10.4}{12.0}\selectfont
    \textbf{Цель работы}\\[0.15cm]
    Сравнить три способа хранения одной и той же модели CRDT-операций:
    реляционный, документный и распределённый.}%
  };

  \node[
    anchor=north west,
    fill=softcyan,
    rounded corners=10pt,
    text width=7.0cm,
    inner xsep=0.35cm,
    inner ysep=0.30cm,
    align=left
  ] at ([xshift=7.60cm,yshift=-1.95cm]current page.north west) {%
    {\fontsize{9.7}{11.3}\selectfont
    \textbf{Что исследуется}\\[0.10cm]
    Один и тот же контракт операции для прототипа коллаборативного календаря
    с оффлайн-режимом и синхронизацией между репликами.}%
  };

  \node[anchor=north west, text=textblue] at ([xshift=1.05cm,yshift=-4.92cm]current page.north west) {%
    {\fontsize{12.2}{13.8}\selectfont\bfseries Основные задачи}%
  };

  \node[anchor=north west, align=left, text width=13.4cm] at ([xshift=1.10cm,yshift=-5.48cm]current page.north west) {%
    {\fontsize{9.5}{11.0}\selectfont
    1. Описать логическую модель операции и структуру репозитория.\\
    2. Спроектировать хранение в PostgreSQL и MongoDB.\\
    3. Рассмотреть распределённый вариант на YTsaurus.\\
    4. Сопоставить ограничения, индексацию и типовые запросы.\\
    5. Сравнить результаты измерений на подготовленных датасетах.}%
  };
\end{tikzpicture}

\footerband{2}
